\scanpage{31}
\deleted{Supposons}
\uline{NB}\quad $G(j),G(j')$ et $N(j)$ sont connus quand on connaît (en plus des \inserted{parties} $G_0(j)$ et $G_0(j')$ de $C_0$) les ensembles
\[
\left\{\begin{aligned}
G^*&=\{x\in C\setminus C_0\mid x\in G(\alpha(x))\}\\
P^*&=\{x\in C\setminus C_0\mid x\in G(\bar\alpha(x))\}
\end{aligned}\right.
\]
où $\bar\alpha(x)$ désigne l'unique $j'\in J$ tel que $j'\ne j=\alpha(x)$.

On dit qu'une position $x\in C\setminus C_0$ qui est $\in G^*\inserted{\text{ (resp. }\in P^*\text{)}}$, i.e. \deleted{telle} \inserted{gagnante} \unsure{pour} le joueur qui a l'\unsure{initiation|initiative} (resp. pour le joueur qui n'a pas l'initiative) est \uline{gagnante} (resp. \uline{perdante}); on dit que c'est une \uline{position nulle} \deleted{autrement} si elle n'est ni gagnante, ni perdante, on pose
\[
N^*=(C\setminus C_0)\setminus(G^*\cup P^*)=N\cap(C\setminus C_0)
\]
(où $N$ est défini dans le corollaire ci-dessus).

On aura
\[
\left\{\begin{aligned}
G(j)\cap C(j)&=G^*\cap C(j)\\
G(j)\cap\bar C(j)&=G(j)\cap C(j')=P^*\cap C(j')\\
G(j)\cap C_0&=G_0(j)
\end{aligned}\right.
\]
ce qui montre que \inserted{les} $G(j)$ $(j\in\unsure{G|J})$ se reconstituent par la connaissance de $G^*,P^*$ et des $G_0(j)$ $(j\in J)$.

Supposons qu'on ait une application
\[
\hat\alpha:C\longrightarrow J
\]
prolongeant $\alpha$, et qu'on ait, pour $\forall x\in C$ et $y\in R(x)$
\[
\hat\alpha(y)\ne\hat\alpha(x)
\]
% archive foliation: 30
